
\subsubsection{6.1 Interpretation of Key Findings}

The results confirm the three-step mechanism at three independent measurement levels. The distributional evidence (h-m1) establishes why the pattern exists. The rank-correlation evidence (h-m2) confirms it without threshold assumptions, using a metric that is robust to score scaling. The AUROC evidence (h-m3) quantifies practical magnitude. This layered structure provides stronger mechanistic support than any single metric.

The approximate symmetry of P1 and P2 magnitudes --- 11.9 AUROC on LLaMA TriviaQA (P1) versus 12.0 AUROC on LLaMA TruthfulQA (P2) --- is consistent with a single underlying distribution-shape driver rather than two separate phenomena. Similarly, the rank-correlation differentials are equal in magnitude but opposite in sign for LLaMA-2-7B ($\Delta\rho = \pm 0.206)$, supporting the interpretation that the same structural asymmetry drives both effects.

Cross-architecture consistency is the strongest available generalizability evidence within the 7B parameter scale. LLaMA-2-7B and Mistral-7B-v0.1 differ in pretraining corpus, architecture details, and tokenization. Directional consistency across both on both benchmark types argues against model-specific confounds.

\subsubsection{6.2 The Raw Sum Finding}

The refutation of P3 --- raw sum is the strongest, not weakest, aggregation on TriviaQA --- is the most practically significant result. The most likely explanation involves length-confidence correlation: on factual recall benchmarks, correct answers tend to be longer, and each token is generated with high confidence. The unnormalized sum then accumulates larger magnitude for correct answers, making it more discriminative than per-token statistics. This interpretation is consistent with raw sum being weakest on TruthfulQA, where all tokens are generated with uniformly high confidence regardless of correctness (making length accumulation non-discriminative).

The h-e1 codebase includes a \texttt{length\textbackslash \_stratified\textbackslash \_auroc} function (validated by unit tests, not yet executed) that would determine whether the raw sum advantage is concentrated in long-answer samples or persists across all length strata. This analysis is necessary to determine whether raw sum represents a genuine sequence-level calibration signal or an artifact of the length distribution in the Farquhar 2024 TriviaQA splits.

\subsubsection{6.3 Comparison to Published Baselines}

The minimum log-probability on TriviaQA (0.849--0.892) exceeds published CCP/mean log-prob AUROC (0.72--0.80) \citep{Fadeeva2024Fact} by 0.05--0.09 AUROC. Raw sum (0.895--0.896) exceeds published Semantic Entropy AUROC ($\approx 0.79)$ \citep{Farquhar2024} by approximately 0.10 AUROC, with a single forward pass compared to SE's 10 forward passes.

These comparisons are cross-pipeline: they span different hardware, dataset subsets, tokenization, and evaluation protocols. The magnitude of the observed gaps --- 0.05--0.10 AUROC --- makes a purely methodological explanation unlikely but cannot be fully excluded without direct within-pipeline comparison. Claims of outperformance over SE should be interpreted as indicating comparable or better performance under the evaluation protocol described here, not as a definitive ranking.

\subsubsection{6.4 Limitations}

\textbf{NQ data gap.} P1 is confirmed on TriviaQA only. NQ inference was executed in h-e1 but the score cache file was not saved. The mechanism prediction for NQ is identical to TriviaQA (same hallucination type: factual recall), and both models confirm P1 with large effect sizes on TriviaQA, but direct NQ confirmation is absent.

\textbf{Single model scale (7B).} All results are at 7B parameters. Lin et al. [2022] report inverse scaling on TruthfulQA, predicting stronger imitative-falsehood effects at larger scales. The 7B P2 result may underestimate the mechanism's magnitude at 70B+.

\textbf{Greedy decoding only.} All results apply within the greedy decoding regime. Generalization to sampling-based settings requires separate validation; greedy log-probabilities may diverge from model uncertainty under temperature > 0.

\textbf{Single peakedness model.} The h-m1 distributional result is established for LLaMA-2-7B only; the corresponding Mistral-7B-v0.1 run encountered a crash. The TruthfulQA peakedness distribution is not directly measured; the flat-distribution characterization of imitative-falsehood hallucinations is an inference consistent with P2 results, not a direct measurement.

\textbf{Multi-sample baselines not in same pipeline.} The comparison to SE and SelfCheckGPT is cross-pipeline as noted above.

\textbf{TruthfulQA label protocol.} Binary labels use ROUGE-L $\geq$ 0.3 following Fadeeva et al. [2024]. Results under alternative labeling schemes (e.g., judge-model protocols) may differ.

\textbf{Unverified citations.} Three 2025 arXiv preprints (Moslonka et al.; Zhang et al.; Ma et al.) are cited but have not been verified against published proceedings. They are included for completeness and should be verified before submission.

\subsubsection{6.5 Broader Implications}

This work provides a principled, label-free selection criterion for zero-cost hallucination detection. The mechanism-grounded taxonomy --- hallucination type determines distribution shape, which determines aggregation optimality --- may generalize to other sequence-level tasks where different error types produce structurally distinct token probability distributions. The finding that aggregation choice matters by up to 12 AUROC points implies that studies using a single aggregation function across diverse benchmarks may be systematically underperforming for specific error classes.

